\StoryParagraph{3}{2026年3月15日，她走后的第七天，助理问我，香港那套房还要不要保留。}
\StoryParagraph{3}{我说，保留。}
\StoryParagraph{3}{他说，没人住，管理费很高。}
\StoryParagraph{3}{我说，我知道。}
\StoryParagraph{3}{他就不说话了。}
\StoryParagraph{3}{那套房里没有她的东西。}
\StoryParagraph{3}{衣服没有，首饰没有，照片没有。她来过几次，每次都只带一个小袋子，离开时也只带走那个小袋子。}
\StoryParagraph{3}{浴室抽屉里倒是留下了三块指甲锉。}
\StoryParagraph{3}{一块粗的，一块细的，还有一块我不知道是做什么的。}
\StoryParagraph{3}{我问助理，这种东西多少钱。}
\StoryParagraph{3}{他说，不贵。}
\StoryParagraph{3}{我说，买三十块。}
\StoryParagraph{3}{他问买来做什么。}
\StoryParagraph{3}{我说，备用。}
\StoryParagraph{3}{第二天，三十块指甲锉送到了房间。助理把它们整齐地放进抽屉，和她留下的三块放在一起。}
\StoryParagraph{3}{这样看起来，像是她从来没有走过。}
\StoryParagraph{3}{又过了几天，仓库管理员给我发来一封邮件，说有人寄回了一份文件。}
\StoryParagraph{3}{四十七页。}
\StoryParagraph{3}{我记得那份文件。最早是四十三页，后来增加了四页。增加的内容我没有看过，是助理说应该补上。}
\StoryParagraph{3}{我问补了什么。}
\StoryParagraph{3}{他说，补了几个账户，几处房产，还有一段时间线。}
\StoryParagraph{3}{我说删掉。}
\StoryParagraph{3}{他说已经装订了。}
\StoryParagraph{3}{文件被放在桌上，封面没有标题，只有一行字：}
\StoryParagraph{3}{“你想知道的，都在里面。”}
\StoryParagraph{3}{这不是我的字。}
\StoryParagraph{3}{我翻开第一页。}
\StoryParagraph{3}{第一页没有她的姓名，只有一个日期。}
\StoryParagraph{3}{2026年3月17日。}
\StoryParagraph{3}{那天我还在等她的电话。}
\StoryParagraph{3}{第二页是一张酒店账单。第三页是一张机票。第四页是一张没有签名的收据。第五页开始，所有内容都变成了我的口吻。}
\StoryParagraph{3}{她说过什么，我做过什么，她什么时候沉默，什么时候看手机，什么时候把帽子摘下来，全部写得很准确。}
\StoryParagraph{3}{准确得像有人一直坐在房间里看着我们。}
\StoryParagraph{3}{我翻到最后一页。}
\StoryParagraph{3}{那里只有一句话：}
\StoryParagraph{3}{“你把我写得太像一个人。”}
\StoryParagraph{3}{我看了很久。}
\StoryParagraph{3}{然后拿起手机，给她发了一条消息。}
\StoryParagraph{3}{我说，文件收到了。}
\StoryParagraph{3}{过了十七分钟，她回复：}
\StoryParagraph{3}{“那不是给你的。”}
\StoryParagraph{3}{我问，那是给谁的。}
\StoryParagraph{3}{她没有回答。}
\StoryParagraph{3}{晚上十一点，仓库管理员又发来一封邮件，说刚刚有人申请把仓库里剩下的东西全部取走。}
\StoryParagraph{3}{我问是谁。}
\StoryParagraph{3}{他说，收件人写的是我的名字。}
\StoryParagraph{3}{我坐在空房间里，听见冰箱启动的声音。}
\StoryParagraph{3}{那一刻我忽然明白，她留下的从来不是东西。}
\StoryParagraph{3}{她留下的是一种方法。}
\StoryParagraph{3}{让一个房间继续等待，让一笔钱继续流动，让一个人不断寻找已经不存在的证据。}
\StoryParagraph{3}{我把三块旧的指甲锉拿出来，试着磨了一下自己的指甲。}
\StoryParagraph{3}{还是尖的。}
\StoryParagraph{3}{我换了一个方向。}
\StoryParagraph{3}{更尖了。}
\StoryParagraph{3}{我没有再磨。}
\StoryParagraph{3}{第二天早上，仓库的东西被送到了办公室。}
\StoryParagraph{3}{两个纸箱，一个硬壳文件袋，还有一只没有寄件人的信封。}
\StoryParagraph{3}{信封很薄，里面只有一张纸。}
\StoryParagraph{3}{不是她的字，也不是我的字。}
\StoryParagraph{3}{“别再找我了。”}
\StoryParagraph{3}{下面没有落款。}
\StoryParagraph{3}{我把纸翻过来，背面是空的。}
\StoryParagraph{3}{我问助理，信是谁送来的。}
\StoryParagraph{3}{他说，前台没见到人。监控里只有一个戴帽子的女人，站在门口三分钟。}
\StoryParagraph{3}{我说，截图给我。}
\StoryParagraph{3}{他说已经调过了。}
\StoryParagraph{3}{照片里的人背对着镜头，帽檐压得很低，手里提着一个纸袋。}
\StoryParagraph{3}{我把照片放大。}
\StoryParagraph{3}{监控的分辨率不够，什么也看不清。}
\StoryParagraph{3}{她的脸像十九年前那张照片一样，逐渐失去细节，只剩下一个轮廓。}
\StoryParagraph{3}{我忽然想起，最早喜欢她的时候，我喜欢的也是一个轮廓。}
\StoryParagraph{3}{那张QQ广告图很小，像素很低。她站在画面中央，脸上有一点笑，旁边有一行我看不清的字。}
\StoryParagraph{3}{我把那张图保存了十九年。}
\StoryParagraph{3}{十九年里，我换过很多手机，很多电脑，很多住处，甚至换过几次身份信息的格式。}
\StoryParagraph{3}{只有那张图没有变。}
\StoryParagraph{3}{它一直那么糊。}
\StoryParagraph{3}{我后来才知道，不是设备不够好。}
\StoryParagraph{3}{是我从来没有真正看清过她。}
\StoryParagraph{3}{2026年3月20日，她离开后的第十二天，我收到一封没有寄件人的邮件。}
\StoryParagraph{3}{邮件说，对方希望我停止寻找，删除所有与她有关的材料，并保留追究责任的权利。}
\StoryParagraph{3}{我问，相关材料包括什么。}
\StoryParagraph{3}{邮件里没有署名，也没有写明对方是谁，只列出所有能够识别出一个人身份的文字、图片、录音和叙述。}
\StoryParagraph{3}{我说，连我自己的回忆也包括吗。}
\StoryParagraph{3}{没有人回答。}
\StoryParagraph{3}{我把邮件关掉。}
\StoryParagraph{3}{那时我还不知道，几个月以后，这句话会被更多人用不同的方式说一遍。}
\StoryParagraph{3}{我让他把办公室里所有打印件都收走。}
\StoryParagraph{3}{他说，已经备份过了。}
\StoryParagraph{3}{我说，全部删掉。}
\StoryParagraph{3}{他说，备份不在我这里。}
\StoryParagraph{3}{我没有再说话。}
\StoryParagraph{3}{下午，助理把排版团队做的几份样稿放在我桌上。}
\StoryParagraph{3}{一份按小说排，一份按情书排，还有一份把里面的数字单独列出来，做成了一张内部校对表。}
\StoryParagraph{3}{他问，要不要在最后加一句解释。}
\StoryParagraph{3}{我问，解释什么。}
\StoryParagraph{3}{他说，解释哪些是故事，哪些不是。}
\StoryParagraph{3}{我说，不加。}
\StoryParagraph{3}{他说，读者会误会。}
\StoryParagraph{3}{我说，误会也是读者的自由。}
\StoryParagraph{3}{助理看了我一会儿，没有再说话。}
\StoryParagraph{3}{那篇文章没有在那天发布。}
\StoryParagraph{3}{我以前也是这样。}
\StoryParagraph{3}{我相信她爱我，也相信她不爱我。}
\StoryParagraph{3}{我相信那四十几页，也相信它一个字都不可信。}
\StoryParagraph{3}{我相信只要钱够多，所有事情都可以重新开始。}
\StoryParagraph{3}{又相信有些事情一旦发生，多少钱也买不回去。}
\StoryParagraph{3}{我对助理说，别再看了。}
\StoryParagraph{3}{他说，我没看。}
\StoryParagraph{3}{我说，你看了。}
\StoryParagraph{3}{他说，我只看了标题。}
\StoryParagraph{3}{我说，标题也是看。}
\StoryParagraph{3}{他点点头，像接受一条新的财务制度。}
\StoryParagraph{3}{那天晚上，我打开电脑，把那篇文章的源文件重新编译了一次。}
\StoryParagraph{3}{我不知道为什么要这么做。}
\StoryParagraph{3}{可能是因为源文件比 PDF 更像事实。}
\StoryParagraph{3}{PDF 已经完成了，文字固定在纸上，不能移动，不能反驳，也不能要求作者重新解释。}
\StoryParagraph{3}{源文件不一样。}
\StoryParagraph{3}{源文件里还有很多空行。}
\StoryParagraph{3}{每一行都可能被删掉，也可能被加上新的内容。}
\StoryParagraph{3}{我在最后一页加了一行注释。}
\StoryParagraph{3}{“以下内容不再更新。”}
\StoryParagraph{3}{保存之前，我又删掉了。}
\StoryParagraph{3}{我不想给它一个结论。}
\StoryParagraph{3}{结论是给别人看的。}
\StoryParagraph{3}{我认识她的第一个月，曾经让人做过一份关于她的报告。}
\StoryParagraph{3}{那时我觉得，了解一个人和收购一家公司差不多。}
\StoryParagraph{3}{先查资产，再查关系，再查风险，最后算一个可以接受的价格。}
\StoryParagraph{3}{报告做得很漂亮。}
\StoryParagraph{3}{封面用的是灰色硬纸，内页是米白色，所有数字右对齐，所有日期统一成国际格式。}
\StoryParagraph{3}{我拿到它的时候，先看了目录。}
\StoryParagraph{3}{项目、合作方、股权、账户、房产、家人、朋友、前任，以及可能的冲突。}
\StoryParagraph{3}{没有“喜欢”这一栏。}
\StoryParagraph{3}{我问负责的人，为什么没有。}
\StoryParagraph{3}{他说，喜欢无法核实。}
\StoryParagraph{3}{我说，喜欢不是事实吗。}
\StoryParagraph{3}{他说，无法核实的事实不能写进报告。}
\StoryParagraph{3}{我说，那就写成推测。}
\StoryParagraph{3}{他说，推测会影响报告可信度。}
\StoryParagraph{3}{我把报告放回桌上。}
\StoryParagraph{3}{那天晚上我第一次意识到，我身边的人很少做没有用的事。}
\StoryParagraph{3}{包括我。}
\StoryParagraph{3}{我对她说过很多次，我相信她。}
\StoryParagraph{3}{但我从来没有把“相信”交给她保管。}
\StoryParagraph{3}{我把它放在账户里，放在文件里，放在助理的日程里，放在每一笔转账的备注里。}
\StoryParagraph{3}{只要它有记录，就算不上真正交出去。}
\StoryParagraph{3}{她知道。}
\StoryParagraph{3}{她只是一直没有说。}
\StoryParagraph{3}{直到那天晚上，她问我，你能不能有一次不问为什么。}
\StoryParagraph{3}{我说，可以。}
\StoryParagraph{3}{她说，那你现在就不要问。}
\StoryParagraph{3}{我说，好。}
\StoryParagraph{3}{后来我还是问了。}
\StoryParagraph{3}{我问她去了哪里，问她什么时候回来，问她还要不要那套房，问她是不是已经决定好了。}
\StoryParagraph{3}{她一个问题也没有回答。}
\StoryParagraph{3}{她只说，你看，你做不到。}
\StoryParagraph{3}{我说，我只是想知道。}
\StoryParagraph{3}{她说，想知道也是一种索取。}
\StoryParagraph{3}{当时我不理解。}
\StoryParagraph{3}{现在也不完全理解。}
\StoryParagraph{3}{但我记住了这句话。}
\StoryParagraph{3}{2026年3月29日，她走后的第三个星期，那个没有寄件人的信封又寄来一次。}
\StoryParagraph{3}{这次里面是一张收据。}
\StoryParagraph{3}{金额很小，只有九十六块。}
\StoryParagraph{3}{购买地点是一家机场里的文具店。}
\StoryParagraph{3}{商品名称写着：黑色记号笔、普通信封、A4打印纸。}
\StoryParagraph{3}{购买时间是她离开后的第二天。}
\StoryParagraph{3}{我拿着收据问助理，这是谁买的。}
\StoryParagraph{3}{他说，查不到。}
\StoryParagraph{3}{我说，机场有监控吗。}
\StoryParagraph{3}{他说，有。}
\StoryParagraph{3}{我说，调出来。}
\StoryParagraph{3}{他说，已经调过了。}
\StoryParagraph{3}{还是那个戴帽子的女人。}
\StoryParagraph{3}{她站在柜台前，先买了纸和信封，然后买了一支黑色记号笔。}
\StoryParagraph{3}{付款时，她把左手放在柜台上。}
\StoryParagraph{3}{我盯着那只手看了很久。}
\StoryParagraph{3}{指甲修得很圆，没有一点锋利的边缘。}
\StoryParagraph{3}{她以前总说，指甲不能刮人。}
\StoryParagraph{3}{我问助理，机场那家店能不能买到指甲锉。}
\StoryParagraph{3}{他说，应该能。}
\StoryParagraph{3}{我说，她为什么没买。}
\StoryParagraph{3}{他说，我不知道。}
\StoryParagraph{3}{我说，你去问。}
\StoryParagraph{3}{他说，问谁。}
\StoryParagraph{3}{我这才发现，我身边有很多人可以查到她去过哪里，却没有一个人能问她为什么没有买一块指甲锉。}
\StoryParagraph{3}{我把收据夹进文件袋里。}
\StoryParagraph{3}{又过了两天，她给我打电话。}
\StoryParagraph{3}{号码没有显示归属地。}
\StoryParagraph{3}{我接起来，没有说话。}
\StoryParagraph{3}{她也没有说话。}
\StoryParagraph{3}{电话里有很轻的风声。}
\StoryParagraph{3}{我问，你在哪里。}
\StoryParagraph{3}{她说，你还是问了。}
\StoryParagraph{3}{我说，对不起。}
\StoryParagraph{3}{她说，不用。}
\StoryParagraph{3}{我问，那你为什么打电话。}
\StoryParagraph{3}{她说，想听听你还会不会接。}
\StoryParagraph{3}{我说，会。}
\StoryParagraph{3}{她说，以后不要了。}
\StoryParagraph{3}{我说，好。}
\StoryParagraph{3}{她说，你每次都说好。}
\StoryParagraph{3}{我说，这次是真的。}
\StoryParagraph{3}{她在那边笑了一下。}
\StoryParagraph{3}{很短。}
\StoryParagraph{3}{我突然想起第一次见她的时候，她也这样笑过。}
\StoryParagraph{3}{不是因为我说了什么。}
\StoryParagraph{3}{是因为服务生把一杯水放错了位置，她把水杯往旁边挪了一点。}
\StoryParagraph{3}{那一刻我觉得她很漂亮。}
\StoryParagraph{3}{现在我想不起来那杯水放在哪里了。}
\StoryParagraph{3}{她问，你把文章删了吗。}
\StoryParagraph{3}{我说，没有。}
\StoryParagraph{3}{她说，为什么。}
\StoryParagraph{3}{我说，已经发出去了。}
\StoryParagraph{3}{她说，你总能找到一个听起来合理的理由。}
\StoryParagraph{3}{我说，那不是理由，是事实。}
\StoryParagraph{3}{她说，你看，又来了。}
\StoryParagraph{3}{我没有说话。}
\StoryParagraph{3}{她问，报告还在吗。}
\StoryParagraph{3}{我说，在。}
\StoryParagraph{3}{她说，烧了吧。}
\StoryParagraph{3}{我说，里面有你的东西。}
\StoryParagraph{3}{她说，里面没有我。}
\StoryParagraph{3}{我说，有你的生日、住址、家人的名字，还有很多——}
\StoryParagraph{3}{她打断我：}
\StoryParagraph{3}{“那是你搜到的，不是我给你的。”}
\StoryParagraph{3}{我说，我只是想了解你。}
\StoryParagraph{3}{她说，你想拥有一个关于我的解释。}
\StoryParagraph{3}{我没有反驳。}
\StoryParagraph{3}{她说，解释比人安全。}
\StoryParagraph{3}{我问，你现在安全了吗。}
\StoryParagraph{3}{她没有回答。}
\StoryParagraph{3}{电话断了。}
\StoryParagraph{3}{我把手机放在桌上，屏幕朝下。}
\StoryParagraph{3}{我已经习惯这样放手机了。}
\StoryParagraph{3}{以前是为了不让她看见消息。}
\StoryParagraph{3}{后来是为了不看见她的消息。}
\StoryParagraph{3}{同一个动作，原因可以完全相反。}
\StoryParagraph{3}{我让人把报告送到楼下。}
\StoryParagraph{3}{地下停车场里有一个金属桶，是清理文件用的。纸张烧起来很快，塑料封皮烧得慢一些，硬纸板会发出很轻的爆裂声。}
\StoryParagraph{3}{助理问，要不要留一份电子版。}
\StoryParagraph{3}{我说，不要。}
\StoryParagraph{3}{他说，万一以后需要。}
\StoryParagraph{3}{我说，不会有以后。}
\StoryParagraph{3}{他说，这句话不能写进合同。}
\StoryParagraph{3}{我看着他。}
\StoryParagraph{3}{他说，我只是提醒。}
\StoryParagraph{3}{他把第一叠纸放进去。}
\StoryParagraph{3}{火从纸边开始，先烧掉的是日期。}
\StoryParagraph{3}{然后是账户。}
\StoryParagraph{3}{然后是她的名字。}
\StoryParagraph{3}{我站在那里，看着一个人被按顺序烧掉。}
\StoryParagraph{3}{我以前以为保存就是珍惜。}
\StoryParagraph{3}{那天才知道，保存也可以是一种暴力。}
\StoryParagraph{3}{最后剩下的是封面。}
\StoryParagraph{3}{灰色的硬纸卷起来，像一个合上的嘴。}
\StoryParagraph{3}{我让助理把它也放进去。}
\StoryParagraph{3}{火灭以后，桶底有一块没有烧干净的纸。}
\StoryParagraph{3}{上面只剩两个字：}
\StoryParagraph{3}{“母亲。”}
\StoryParagraph{3}{这是报告里唯一不是我让人写进去的词。}
\StoryParagraph{3}{我不知道它从哪里来。}
\StoryParagraph{3}{也许是助理整理材料时弄错了。}
\StoryParagraph{3}{也许是我自己加的，只是忘了。}
\StoryParagraph{3}{我把那块纸捡起来，放进钱包。}
\StoryParagraph{3}{后来它在钱包里待了三个多月。}
\StoryParagraph{3}{我每次付款都会摸到它。}
\StoryParagraph{3}{像摸到一块很小的伤口。}
\StoryParagraph{3}{2026年4月17日，她走后的第四十天，我去了那家电影院。}
\StoryParagraph{3}{不是为了看电影。}
\StoryParagraph{3}{我想知道二十六个人被请出去以后，影院还剩下什么。}
\StoryParagraph{3}{经理认出了我。}
\StoryParagraph{3}{他说，您那天包场以后，网上有很多人来问。}
\StoryParagraph{3}{我说，问什么。}
\StoryParagraph{3}{他说，问那场电影是不是只有两个人。}
\StoryParagraph{3}{我说，是。}
\StoryParagraph{3}{他说，他们都觉得很浪漫。}
\StoryParagraph{3}{我说，不浪漫。}
\StoryParagraph{3}{他说，为什么。}
\StoryParagraph{3}{我说，因为当时她一直在看手机。}
\StoryParagraph{3}{经理不知道该怎么接。}
\StoryParagraph{3}{我买了两张最普通的票。}
\StoryParagraph{3}{这次没有清场。}
\StoryParagraph{3}{影厅里有十几个人，最后一排坐着一对年轻情侣，前面有一个抱着爆米花的小孩。}
\StoryParagraph{3}{我坐在中间。}
\StoryParagraph{3}{电影开始后，那个小孩一直问旁边的大人，为什么兔子不能开车，为什么狐狸不说实话，为什么坏人最后又变好了。}
\StoryParagraph{3}{大人不停地让他小声一点。}
\StoryParagraph{3}{我却觉得这个声音很好。}
\StoryParagraph{3}{有人在我前面挡住了字幕。}
\StoryParagraph{3}{有人把饮料碰倒了。}
\StoryParagraph{3}{有人在最安静的地方咳嗽。}
\StoryParagraph{3}{电影结束时，没有人鼓掌。}
\StoryParagraph{3}{灯亮起来，大家一起站起来，慢慢走出去。}
\StoryParagraph{3}{没有人多看我一眼。}
\StoryParagraph{3}{我在空掉的座位上坐了五分钟。}
\StoryParagraph{3}{那天我终于知道，别人不认识你，不代表世界是空的。}
\StoryParagraph{3}{只是世界没有为你暂停。}
\StoryParagraph{3}{回到家，门口放着一只纸箱。}
\StoryParagraph{3}{寄件人写的是那家电影院。}
\StoryParagraph{3}{里面是一件孩子的外套，深蓝色，袖口有一块洗不掉的爆米花油渍。}
\StoryParagraph{3}{经理说，那个小孩走的时候把外套落下了，工作人员以为是我的。}
\StoryParagraph{3}{我说，不是我的。}
\StoryParagraph{3}{经理说，我知道，但地址只有这里。}
\StoryParagraph{3}{我让助理送去失物招领。}
\StoryParagraph{3}{助理说，已经问过了，没有人登记。}
\StoryParagraph{3}{我说，那就捐掉。}
\StoryParagraph{3}{他说，好的。}
\StoryParagraph{3}{我看着那件外套，忽然想起她说过想要一个孩子。}
\StoryParagraph{3}{我没有问她为什么一定要。}
\StoryParagraph{3}{她也没有告诉我。}
\StoryParagraph{3}{我以前把所有没有答案的问题都归到钱上。}
\StoryParagraph{3}{因为钱可以解释很多事情。}
\StoryParagraph{3}{但不能解释一个人为什么想成为母亲，也不能解释一个人为什么不愿意成为你故事里的母亲。}
\StoryParagraph{3}{那天晚上，我第一次打开了她的公开主页。}
\StoryParagraph{3}{里面有很多工作照片。}
\StoryParagraph{3}{她站在海边，站在片场，站在灯光下，站在人群中间。每一张照片都有很多人，只有她看起来像一个人。}
\StoryParagraph{3}{我往前翻。}
\StoryParagraph{3}{翻到很久以前的一张。}
\StoryParagraph{3}{她穿着一件白色上衣，头发比现在长，脸还没有被高清相机保存下来。}
\StoryParagraph{3}{那张照片下面有很多评论。}
\StoryParagraph{3}{有人说她漂亮。}
\StoryParagraph{3}{有人说她像某个演员。}
\StoryParagraph{3}{有人问她什么时候进组。}
\StoryParagraph{3}{没有人问她今天过得好不好。}
\StoryParagraph{3}{我关掉了页面。}
\StoryParagraph{3}{又打开。}
\StoryParagraph{3}{再关掉。}
\StoryParagraph{3}{最后我把那个网页从浏览器历史记录里删了。}
\StoryParagraph{3}{这件事没有任何意义。}
\StoryParagraph{3}{但我还是做了。}
\StoryParagraph{3}{人有时候需要一些没有意义的动作，证明自己还能够决定什么。}
\StoryParagraph{3}{2026年5月7日，她离开后的第六十天，我收到一条短信。}
\StoryParagraph{3}{不是她的号码。}
\StoryParagraph{3}{内容只有一句：}
\StoryParagraph{3}{“你还留着那张旧照片吗？”}
\StoryParagraph{3}{我回复：}
\StoryParagraph{3}{“留着。”}
\StoryParagraph{3}{对方问：}
\StoryParagraph{3}{“为什么？”}
\StoryParagraph{3}{我打了很多字。}
\StoryParagraph{3}{删掉。}
\StoryParagraph{3}{最后只回复了两个字：}
\StoryParagraph{3}{“忘了。”}
\StoryParagraph{3}{过了很久，对方又发来一句：}
\StoryParagraph{3}{“那就继续忘。”}
\StoryParagraph{3}{我知道是她。}
\StoryParagraph{3}{但我没有问。}
\StoryParagraph{3}{第二天，我把照片从旧手机里导出来。}
\StoryParagraph{3}{这次没有传到新设备。}
\StoryParagraph{3}{我把它存进一个没有名字的文件夹，又给文件夹加了密码。}
\StoryParagraph{3}{密码不是她的生日，也不是我的生日。}
\StoryParagraph{3}{是第一次见她那天，服务生把水杯放错位置的时间。}
\StoryParagraph{3}{我记得。}
\StoryParagraph{3}{下午三点十七分。}
\StoryParagraph{3}{我不知道为什么记得。}
\StoryParagraph{3}{可能因为那是少数一件没有被记录在报告里的事。}
\StoryParagraph{3}{2026年5月22日，网上出现了一则匿名消息。}
\StoryParagraph{3}{消息里没有姓名，只说一个J姓女演员和一个九零后富豪，涉及海外生育安排和巨额资金纠纷。}
\StoryParagraph{3}{很多人立刻开始猜她是谁。}
\StoryParagraph{3}{她的工作室很快发了声明，否认相关传言，说已经取证，并会通过法律途径处理。}
\StoryParagraph{3}{我看到了那份声明。}
\StoryParagraph{3}{我没有转发，也没有评论。}
\StoryParagraph{3}{那天晚上，我把文章的文件夹改了名字。}
\StoryParagraph{3}{从“她”改成“草稿”。}
\StoryParagraph{3}{2026年6月3日，我把那四十几页材料交给律师看。}
\StoryParagraph{3}{他说，这些只能说明我保存过它们，不能说明里面每一件事都发生过。}
\StoryParagraph{3}{我问，钱呢。}
\StoryParagraph{3}{他说，钱要看转账凭证、用途和双方的约定。}
\StoryParagraph{3}{我问，口头说过的算不算。}
\StoryParagraph{3}{他说，要看谁说、什么时候说、有没有别的证据。}
\StoryParagraph{3}{我问，那这篇文章算不算。}
\StoryParagraph{3}{他说，文章首先是文章。}
\StoryParagraph{3}{我说，我知道。}
\StoryParagraph{3}{他说，网上那则匿名消息和这份材料没有直接关系。}
\StoryParagraph{3}{我问，对方有没有提出新的要求。}
\StoryParagraph{3}{他说，没有。}
\StoryParagraph{3}{我又问，她有没有问起我。}
\StoryParagraph{3}{他说，这不属于他的工作范围。}
\StoryParagraph{3}{律师走后，我让助理把所有日程清空。}
\StoryParagraph{3}{包括预留的房间、未使用的机票、没有住过的酒店、没有戴过的戒指，还有一顿已经付过定金的晚餐。}
\StoryParagraph{3}{助理问，晚餐取消吗。}
\StoryParagraph{3}{我说，取消。}
\StoryParagraph{3}{他说，定金不能退。}
\StoryParagraph{3}{我说，认赔。}
\StoryParagraph{3}{他说，好的。}
\StoryParagraph{3}{他拿出手机，一个一个打电话。}
\StoryParagraph{3}{我坐在窗边，看着他把一顿不存在的晚餐从世界上删掉。}
\StoryParagraph{3}{他打到最后一家餐厅时，对方说可以把定金改成一张长期储值卡。}
\StoryParagraph{3}{助理问我行不行。}
\StoryParagraph{3}{我说，不要。}
\StoryParagraph{3}{他说，这样损失会小很多。}
\StoryParagraph{3}{我说，就让它损失。}
\StoryParagraph{3}{他说，明白了。}
\StoryParagraph{3}{我知道他不明白。}
\StoryParagraph{3}{很多时候，只有损失真正发生，事情才算结束。}
\StoryParagraph{3}{否则它只是换了一个账户继续存在。}
\StoryParagraph{3}{2026年6月8日，三个月以后，我把香港的房子卖掉了。}
\StoryParagraph{3}{买家是一对夫妻。}
\StoryParagraph{3}{男的看房时一直在计算面积，女的站在窗边看海。}
\StoryParagraph{3}{女的问，这里的夕阳是不是很好看。}
\StoryParagraph{3}{我说，我没看过。}
\StoryParagraph{3}{她回头看我，像是不相信。}
\StoryParagraph{3}{我说，窗帘一直拉着。}
\StoryParagraph{3}{她问，为什么。}
\StoryParagraph{3}{我说，她睡觉不能有光。}
\StoryParagraph{3}{女的点点头，没有再问。}
\StoryParagraph{3}{签合同的时候，买家的孩子在客厅里跑来跑去。}
\StoryParagraph{3}{他撞到了墙角，哭了两声，很快又跑开了。}
\StoryParagraph{3}{我站在原来的位置，看着这个房间第一次真正有人生活。}
\StoryParagraph{3}{他们没有把所有家具换掉。}
\StoryParagraph{3}{他们留下了冰箱，留下了客厅的灯，也留下了浴室里那面太大的镜子。}
\StoryParagraph{3}{临走时，那个女人问我，您为什么卖掉这里。}
\StoryParagraph{3}{我说，房子太大。}
\StoryParagraph{3}{她说，两个人住正好。}
\StoryParagraph{3}{我说，我们不是两个人。}
\StoryParagraph{3}{她没有听懂。}
\StoryParagraph{3}{我也没有解释。}
\StoryParagraph{3}{那天晚上，我住进了一间只有四十平方米的公寓。}
\StoryParagraph{3}{没有海景，没有助理，没有单独的衣帽间。}
\StoryParagraph{3}{卧室和客厅之间只有一面半高的墙，冰箱启动时，床也会轻轻震一下。}
\StoryParagraph{3}{我第一次在半夜听见楼上的脚步声。}
\StoryParagraph{3}{有人拖椅子，有人冲水，有人把玻璃瓶放在地上。}
\StoryParagraph{3}{这些声音以前会让她睡不着。}
\StoryParagraph{3}{我以前会让整栋楼安静。}
\StoryParagraph{3}{现在我没有办法。}
\StoryParagraph{3}{奇怪的是，我渐渐睡着了。}
\StoryParagraph{3}{搬家后的第五天，我在门缝下面发现一张明信片。}
\StoryParagraph{3}{没有邮戳，只有一幅很普通的海。}
\StoryParagraph{3}{海水是蓝色的，天空也是蓝色的，远处有一艘白船。}
\StoryParagraph{3}{背面写着：}
\StoryParagraph{3}{“这里晚上有声音。”}
\StoryParagraph{3}{没有署名。}
\StoryParagraph{3}{我把明信片放在桌上。}
\StoryParagraph{3}{晚上，楼上的孩子哭了很久。}
\StoryParagraph{3}{我没有让人去敲门。}
\StoryParagraph{3}{第二天早上，哭声停了。}
\StoryParagraph{3}{我出门时，楼道里有一辆小推车，里面放着牛奶、纸尿裤和一袋橘子。}
\StoryParagraph{3}{我站在那里看了一会儿。}
\StoryParagraph{3}{邻居开门出来，是一个头发很乱的年轻女人。}
\StoryParagraph{3}{她对我说，对不起，昨天晚上吵到你了。}
\StoryParagraph{3}{我说，没有。}
\StoryParagraph{3}{她说，孩子刚出生，不知道什么时候会哭。}
\StoryParagraph{3}{我说，正常。}
\StoryParagraph{3}{她愣了一下，像是没想到我会这么说。}
\StoryParagraph{3}{我问，男孩还是女孩。}
\StoryParagraph{3}{她说，女孩。}
\StoryParagraph{3}{我说，恭喜。}
\StoryParagraph{3}{她笑了，说谢谢。}
\StoryParagraph{3}{我下楼买咖啡。}
\StoryParagraph{3}{回来时，她把门关上了。}
\StoryParagraph{3}{楼道里只剩下婴儿很轻的声音。}
\StoryParagraph{3}{我站在门口，突然想起那条被我划掉的提醒：预产期前三个月，准备房间。}
\StoryParagraph{3}{我那时只用了半秒。}
\StoryParagraph{3}{手指向左一划，提醒就消失了。}
\StoryParagraph{3}{我一直以为自己删除的是一个日期。}
\StoryParagraph{3}{后来才知道，我删除的是一个还没有出生的人。}
\StoryParagraph{3}{2026年7月6日，她离开后的第一百二十天，我又接到那个陌生号码的电话。}
\StoryParagraph{3}{她说，最近好吗。}
\StoryParagraph{3}{我说，还好。}
\StoryParagraph{3}{她说，你住小房子了。}
\StoryParagraph{3}{我说，你怎么知道。}
\StoryParagraph{3}{她说，我只是猜的。}
\StoryParagraph{3}{我说，猜得很准。}
\StoryParagraph{3}{她说，终于有一次不是你买的。}
\StoryParagraph{3}{我说，是租的。}
\StoryParagraph{3}{她笑了一下。}
\StoryParagraph{3}{我问，你在哪里。}
\StoryParagraph{3}{她说，不能问。}
\StoryParagraph{3}{我说，好。}
\StoryParagraph{3}{她说，你真的变了。}
\StoryParagraph{3}{我说，哪里。}
\StoryParagraph{3}{她说，以前你会让人查。}
\StoryParagraph{3}{我说，现在也可以。}
\StoryParagraph{3}{她说，那你为什么不查。}
\StoryParagraph{3}{我说，因为你说不能问。}
\StoryParagraph{3}{电话那边安静了一会儿。}
\StoryParagraph{3}{她说，你以前也听我的话。}
\StoryParagraph{3}{我说，以前不是听，是执行。}
\StoryParagraph{3}{她没有笑。}
\StoryParagraph{3}{我也没有。}
\StoryParagraph{3}{她问，房子卖了多少钱。}
\StoryParagraph{3}{我说，不重要。}
\StoryParagraph{3}{她说，你终于学会说这个了。}
\StoryParagraph{3}{我说，可能吧。}
\StoryParagraph{3}{她说，那张照片呢。}
\StoryParagraph{3}{我说，还在。}
\StoryParagraph{3}{她说，你应该删掉。}
\StoryParagraph{3}{我说，删了就没有了。}
\StoryParagraph{3}{她说，本来就没有。}
\StoryParagraph{3}{我没有说话。}
\StoryParagraph{3}{她说，小时候我最讨厌拍照。}
\StoryParagraph{3}{我问，为什么。}
\StoryParagraph{3}{她说，拍完之后大家都说，和本人不像。}
\StoryParagraph{3}{我问，哪里不像。}
\StoryParagraph{3}{她说，照片里的人不会动，也不会说错话。}
\StoryParagraph{3}{我说，这不是很多人喜欢照片的原因吗。}
\StoryParagraph{3}{她说，所以我不喜欢。}
\StoryParagraph{3}{这是她第一次主动告诉我一件小时候的事。}
\StoryParagraph{3}{我没有追问。}
\StoryParagraph{3}{她说，你学会了。}
\StoryParagraph{3}{我说，学会什么。}
\StoryParagraph{3}{她说，留一点空白。}
\StoryParagraph{3}{电话又安静下来。}
\StoryParagraph{3}{这一次，是我先挂的。}
\StoryParagraph{3}{我把手机放在桌上，屏幕朝上。}
\StoryParagraph{3}{屏幕亮了很久，没有新的消息。}
\StoryParagraph{3}{我没有把它扣过去。}
\StoryParagraph{3}{2026年8月27日，第一篇文章在网上发布以后，我重新打开源文件。}
\StoryParagraph{3}{那天晚上，她的工作室也发了声明。}
\StoryParagraph{3}{声明没有讨论文章里的细节，只说艺人的声誉不能被拿来交换，并表示一切交给法院处理，后续只公布维权进展。}
\StoryParagraph{3}{在原来的结尾下面，我增加了一页。}
\StoryParagraph{3}{这一页没有故事。}
\StoryParagraph{3}{没有金额，没有地点，没有酒店，没有飞机，没有账户，也没有她的名字。}
\StoryParagraph{3}{我只写了几行：}
\StoryParagraph{3}{“她问我最近好吗。}
\StoryParagraph{3}{我说，还好。}
\StoryParagraph{3}{这一次，我没有告诉她我在哪里。}
\StoryParagraph{3}{她也没有继续问。”}
\StoryParagraph{3}{我编译了两次。}
\StoryParagraph{3}{第一次，系统提示有一处引用没有找到。}
\StoryParagraph{3}{第二次，没有错误。}
\StoryParagraph{3}{原来的 PDF 一共二十二页，新的 PDF 是二十三页。}
\StoryParagraph{3}{比第一篇多一页。}
\StoryParagraph{3}{我把文件命名为 \\texttt{reply.pdf}。}
\StoryParagraph{3}{没有上传。}
\StoryParagraph{3}{2026年8月28日凌晨，助理问，第二篇什么时候发布。}
\StoryParagraph{3}{我问，谁告诉你有第二篇。}
\StoryParagraph{3}{他说，大家都在等。}
\StoryParagraph{3}{我说，等什么。}
\StoryParagraph{3}{他说，等她回头，等你反转，等你公布证据，等你把故事写完。}
\StoryParagraph{3}{我说，故事已经完了。}
\StoryParagraph{3}{他说，结尾不是开放的吗。}
\StoryParagraph{3}{我说，开放不是邀请。}
\StoryParagraph{3}{他说，明白。}
\StoryParagraph{3}{他停了一下，又问，那文件要不要备份。}
\StoryParagraph{3}{我说，备份一份。}
\StoryParagraph{3}{他说，放在哪里。}
\StoryParagraph{3}{我说，放在我看不到的地方。}
\StoryParagraph{3}{他说，这种地方不存在。}
\StoryParagraph{3}{我说，那就放在一个我不会主动寻找的地方。}
\StoryParagraph{3}{他点点头。}
\StoryParagraph{3}{2026年8月28日早上，她在公开账号上发了一段话。}
\StoryParagraph{3}{她说，她不会拿爱情换钱，也不会拿灵魂换钱。她相信法律，相信最后会有公平。}
\StoryParagraph{3}{我看完以后，把手机放在桌上。}
\StoryParagraph{3}{上午，律师来电话。}
\StoryParagraph{3}{他说，案件已经立案，原告申请了财产保全，但对方提出了管辖权异议。}
\StoryParagraph{3}{我问，什么时候开庭。}
\StoryParagraph{3}{他说，现在还不知道，案件尚未进入实体审理。}
\StoryParagraph{3}{我问，三千余万元的争议，是不是只有这三千余万元。}
\StoryParagraph{3}{他说，法院会看正式材料，不会看小说里的数字。}
\StoryParagraph{3}{我说，小说里的数字呢。}
\StoryParagraph{3}{他说，不属于他的工作范围。}
\StoryParagraph{3}{我挂了电话。}
\StoryParagraph{3}{那天我第一次觉得，一篇文章和一场诉讼之间，距离不是纸张的厚度。}
\StoryParagraph{3}{是一个人能不能把自己写进去，另一个人能不能把自己从里面拿出来。}
\StoryParagraph{3}{我把那张旧照片打印出来。}
\StoryParagraph{3}{打印店的机器比十九年前清楚很多，但照片没有变得更清楚。}
\StoryParagraph{3}{她的脸依然模糊。}
\StoryParagraph{3}{我把照片裁掉，只留下背景里那扇窗。}
\StoryParagraph{3}{窗外有一片很亮的白。}
\StoryParagraph{3}{可能是灯，也可能是天空。}
\StoryParagraph{3}{我不知道。}
\StoryParagraph{3}{我把这张照片放进抽屉，没有加密码。}
\StoryParagraph{3}{抽屉里还有三十块新的指甲锉。}
\StoryParagraph{3}{我拿出一块，沿着同一个方向磨。}
\StoryParagraph{3}{这一次没有磨到发亮。}
\StoryParagraph{3}{刚好不刮人就停了。}
\StoryParagraph{3}{我把它放回去。}
\StoryParagraph{3}{晚上，楼上的小女孩又哭了。}
\StoryParagraph{3}{哭声从天花板传下来，很轻，断断续续，像有人在很远的地方敲门。}
\StoryParagraph{3}{我躺在床上，没有叫任何人去处理。}
\StoryParagraph{3}{窗外有车经过，轮胎压过积水，声音很长。}
\StoryParagraph{3}{我想起她说过，海很大。}
\StoryParagraph{3}{那时候我以为她在说海。}
\StoryParagraph{3}{后来我才知道，她说的是一间房子，是一个人，是我们之间那些没有说出口的东西。}
\StoryParagraph{3}{它们都很大。}
\StoryParagraph{3}{大到一个人站在里面，会听见自己的回声。}
\StoryParagraph{3}{凌晨三点，我起床倒水。}
\StoryParagraph{3}{手机亮了一下。}
\StoryParagraph{3}{是一条新消息。}
\StoryParagraph{3}{“你还会写第三篇吗？”}
\StoryParagraph{3}{我看着这句话，没有回复。}
\StoryParagraph{3}{过了一会儿，屏幕自动熄灭。}
\StoryParagraph{3}{房间重新安静下来。}
\StoryParagraph{3}{我把杯子放在桌上，走到窗边。}
\StoryParagraph{3}{楼下有一盏路灯坏了，亮一下，灭一下。}
\StoryParagraph{3}{远处有人开门，婴儿又哭了。}
\StoryParagraph{3}{这一次，哭声没有停。}
\StoryParagraph{3}{我站在那里听了很久。}
\StoryParagraph{3}{直到天快亮的时候，我才回到床上。}
\StoryParagraph{3}{手机还在桌上。}
\StoryParagraph{3}{那条消息没有被删除。}
\StoryParagraph{3}{我也没有回复。}
\StoryParagraph{3}{窗帘没有拉上。}
\StoryParagraph{3}{天亮以后，房间里会有光。}
\StoryParagraph{3}{我想，这没有关系。}
\StoryParagraph{3}{我已经学会了睡在里面。}
